\documentclass[11pt]{article}
\usepackage[margin=1in]{geometry}
\usepackage{amsmath,amssymb,amsfonts,mathtools,bm}
\usepackage[T1]{fontenc}
\usepackage[utf8]{inputenc}
\title{Inline and Display Mathematics}
\author{Source-backed grouped formula sample}
\date{}
\begin{document}
\maketitle
\section{Expressions}
\paragraph{Expression 1.} The following expression is taken from a source-backed formula corpus.

\begin{align*}\left(u^{~\underline{a}}_{\underline{m}}\right) =\left(u^{a}_{\underline{m}}, u^{~i}_{\underline{m}}\right) ~~ \in ~~ SO(1,D-1)~~~~\Leftrightarrow ~~~~u^{~\underline{a}}_{\underline{m}} \eta^{\underline{m}\underline{n}}u^{~\underline{b}}_{\underline{n}} = \eta^{\underline{a}\underline{b}}\end{align*}

\paragraph{Expression 2.} The following expression is taken from a source-backed formula corpus.

\begin{align*}{\cal A}={\cal A}_0+{\cal A}^{\rm int},~~~~~~{\cal A}_0=\int_{t_a}^{t_b} dt\left[\frac{M}{2}\dot x^2-\frac{ \Omega ^2}{2}x^2\right] ,~~~~{\cal A}^{\rm int}=\int_{t_a}^{t_b} dt\,V^{\rm int}_{\Omega}\equiv\int_{t_a}^{t_b} dt\left[-\frac{( \omega ^2- \Omega ^2)}{2}x^2-V^{\rm int}(x)\right] ,\end{align*}

\paragraph{Expression 3.} The following expression is taken from a source-backed formula corpus.

\begin{align*}{\bf Z({\bf x})}=\frac{1}{{\bf x}}+\sum_{q\neq 0}\{\frac{1}{{\bf x-q}}+\frac{1}{{\bf q}}+\frac{1}{{\bf q}}{\bf x}\frac{1}{{\bf q}}+\frac{1}{{\bf q}}{\bf x}\frac{1}{{\bf q}}{\bf x}\frac{1}{{\bf q}}+\frac{1}{{\bf q}}{\bf x}\frac{1}{{\bf q}}{\bf x}\frac{1}{{\bf q}}{\bf x}\frac{1}{{\bf q}}\},\end{align*}

\paragraph{Expression 4.} The following expression is taken from a source-backed formula corpus.

\begin{align*}1 = \int d^{2} {\bf r} \, | \Psi({\bf r})|^2 =A_{l}^{2} \, 2\pi \kappa^{-2} \, \left\{{\mathcal K}_{i \Theta} (\kappa a) + \left( \frac{\kappa }{\tilde{k} } \right)^{2}\,\left[\frac{K_{i \Theta} (\kappa a) }{J_{l} ( \tilde{k} a ) } \right]^{2}\,{\mathcal J_{l}}( \tilde{k} a ) \right\}\; ,\end{align*}

\paragraph{Expression 5.} The following expression is taken from a source-backed formula corpus.

\begin{align*}\dot{\epsilon}_{\delta\sigma}=\epsilon_{\delta\sigma}\left[3{\dot{\bar{a}}_R\over\bar{a}_R}(1+w)+2{\dot{\bar{c}}_4\over\bar{c}_4}{{1+\epsilon_{\delta\sigma}}\over{1+\epsilon}}+\left({\dot{\bar{G}}_4\over\bar{G}_4}-4{\dot{\bar{c}}_4\over\bar{c}_4}\right)(1+\epsilon_{\delta\sigma})\right].\end{align*}
\end{document}
